\section{总结与延伸}

\subsection{核心要点回顾}

本节课围绕 Score-Based Models 展开，系统地建立了以下核心概念：

\begin{enumerate}
    \item \textbf{Score Function}：定义为 $\nabla_x \log p(x)$，它是一个向量场，指向概率增大的方向，完全编码了概率分布的信息。

    \item \textbf{噪声预测即 Score 学习}：DDPM 中的噪声预测器 $\epsilon_\theta$ 与 score function 之间只差一个缩放因子，揭示了噪声预测的深层物理意义。

    \item \textbf{Tweedie's Formula}：提供了从 score function 到后验均值估计（去噪）的直接公式：$\mathbb{E}[\mu \mid x] = x + \sigma^2 \nabla_x \log p(x)$。

    \item \textbf{Langevin Dynamics}：一种仅依赖 score function 的采样方法，通过``梯度上升 + 噪声注入''的迭代过程收敛到目标分布。

    \item \textbf{Denoising Score Matching}：将不可计算的 score matching 目标转化为可计算的去噪目标，其训练过程与 DDPM 完全等价。

    \item \textbf{多尺度噪声与退火策略}：通过多个噪声级别的 score function 和退火 Langevin dynamics 解决低密度区域 score 不准确的问题，这与 DDPM 的逆向去噪过程等价。
\end{enumerate}

\subsection{拓展阅读}

\begin{itemize}
    \item \textbf{Song \& Ermon (2019)}：\textit{Generative Modeling by Estimating Gradients of the Data Distribution} -- 提出 NCSN，首次将 score matching 与退火 Langevin dynamics 结合用于图像生成。
    \item \textbf{Ho et al. (2020)}：\textit{Denoising Diffusion Probabilistic Models (DDPM)} -- 从变分推断角度推导扩散模型训练目标。
    \item \textbf{Song et al. (2021)}：\textit{Score-Based Generative Modeling through Stochastic Differential Equations} -- 提出 Score SDE 框架，将 DDPM 和 NCSN 统一为连续时间过程。
    \item \textbf{Vincent (2011)}：\textit{A Connection Between Score Matching and Denoising Autoencoders} -- 首次建立 denoising score matching 的理论基础。
    \item \textbf{Efron (2011)}：\textit{Tweedie's Formula and Selection Bias} -- Tweedie's formula 在统计学中的经典论述。
\end{itemize}

%% --- Fin del contenido --- %%

\end{document}
